\begin{tikzpicture}[font=\small,line cap=round,line join=round]\draw[geoTrait](0.0,0)rectangle(0.9,.9);\node[font=\small]at(0.45,.45){1};\draw[geoTrait](1.1,0)rectangle(2.0,.9);\node[font=\small]at(1.55,.45){2};\draw[geoTrait](2.2,0)rectangle(3.1,.9);\node[font=\small]at(2.6500000000000004,.45){3};\draw[geoTrait](3.3000000000000003,0)rectangle(4.2,.9);\node[font=\small]at(3.7500000000000004,.45){4};\draw[geoTrait](4.4,0)rectangle(5.300000000000001,.9);\node[font=\small]at(4.8500000000000005,.45){5};\draw[geoTrait](5.5,0)rectangle(6.4,.9);\node[font=\small]at(5.95,.45){6};\geoSolution{\draw[geoReponse](1.1,0)rectangle(2.0,.9);\draw[geoReponse](3.3000000000000003,0)rectangle(4.2,.9);\draw[geoReponse](5.5,0)rectangle(6.4,.9);\node[right,font=\tiny,otGreenDark]at(6.7,.45){trois issues paires};}\end{tikzpicture}\par\vspace{1.5mm}{\otRegular\fontsize{8}{10}\selectfont\color{otMuted}Les six issues du dé ; celles qui réalisent « obtenir un nombre pair » sont entourées en vert.\par}
